\section{Implications}
\label{sec:implications}

\paragraph{Report the harness.} If $\textit{Agent} = \textit{Harness}(\textit{Model})$, then a benchmark score belongs to the pair, not to the model. %EVIDENCE%
We suggest that agent results report the harness alongside the model, described primitive by primitive: the tools offered and their output formats (\prim{Output}), how context is managed (\prim{Context}), the stopping rules and any fixed pipeline (\prim{Control flow}), and the decoding settings (\prim{Model interface}). Comparisons between models are then fair only within one harness, and comparisons between harnesses only with one model.

\paragraph{Locate failures.} When an agent fails, the primitives give a short list of places to look. The model may never have been shown what it needed (\prim{Context}); it may have asked for the right action and the action may not have happened as intended (\prim{Output}); the loop may have stopped too early or run too long (\prim{Control flow}); the response may have been misread (\prim{Model interface}). Only when all of these are ruled out does the failure belong to the model. In our own work, building the harness around a capable model first, until it passed every task in our evaluation, meant that every failure that remained when a smaller model was swapped in could be attributed to the model.

\paragraph{Port by changing one primitive.} Because only \prim{Model interface} knows which model is inside, a harness can be moved to a new model or provider by changing that primitive alone, translating requests and responses to the new format. Harnesses that let model-specific details leak into other primitives (assuming a provider's tool-call format in \prim{Output}, or a provider's error messages in \prim{Context} management) lose this property.

\paragraph{Train the model in its harness.} A model post-trained for many harnesses must guess which one it is in; a model trained inside one harness learns its tool format, the shape of its results and its stopping rule. The primitives make explicit what such training must match: the context layout the model will see (\prim{Context}), the format in which it must request actions and the results it will read (\prim{Output}), and the loop it will run in (\prim{Control flow}). In our experiments, a small model instruction-tuned on verified sessions recorded in its own harness rose from 2 to 18 of 32 held-out tasks in that harness, while the same base model post-trained by its developers for general use scored 14.
